% =====================================================================
%  preamble/tables.tex
%  Sistema de tablas SIMPLIFICADO.
%  Un único entorno DatosTabla basado en longtable + booktabs, con:
%    - título opcional (negrita, centrado, SIN numeración de "Cuadro N")
%    - líneas horizontales de separación entre filas (\hline en cada fila)
%
%  Uso:
%    \begin{DatosTabla}{cols}{encabezado}[título]
%      fila1 & ... \\ \hline
%      fila2 & ... \\ \hline
%      última & ... \\            % la última sin \hline
%    \end{DatosTabla}
%
%  Donde:
%    - {cols}       especificación de columnas (L{}, C{}, R{})
%    - {encabezado} celdas del encabezado separadas por & (sin \\ final)
%    - [título]     OPCIONAL. Texto del título (sin "Cuadro N:")
%  Cada fila del cuerpo termina en "\\ \hline" para la línea de separación.
% =====================================================================

% --- Tipos de columna reutilizables (ancho fijo, suma = ancho de tabla) ---
\newcolumntype{L}[1]{>{\raggedright\arraybackslash}p{#1}}                 % ancho fijo, izq.
\newcolumntype{C}[1]{>{\centering\arraybackslash}p{#1}}                   % ancho fijo, centrado
\newcolumntype{R}[1]{>{\raggedleft\arraybackslash}p{#1}}                  % ancho fijo, der.

% --- Título de tabla: negrita, centrado, sin numeración (no es "Cuadro N") ---
\newcommand{\TituloTabla}[1]{\caption*{\textbf{#1}}}

% --- Ancho del caption de longtable = ancho de texto (evita overfull) ---
\setlength{\LTcapwidth}{\textwidth}

% --- Comando para identificadores largos (parte en cualquier carácter) ---
\newcommand{\TablaClave}[1]{{\ttfamily\seqsplit{#1}}}

% --- Entorno de tabla de datos con cabecera repetible en cada página ---
%   Título opcional (#3). Las filas del cuerpo llevan "\hline" al final
%   para dibujar la línea de separación horizontal entre filas.
\NewDocumentEnvironment{DatosTabla}{mmo}
{%
  \setlength{\tabcolsep}{5pt}%
  \renewcommand{\arraystretch}{1.3}%
  \setlength{\LTleft}{0pt}%
  \setlength{\LTright}{0pt}%
  \begin{longtable}{#1}%
  \IfValueT{#3}{\TituloTabla{#3}\\[3pt]}
    \toprule
    #2 \\ \midrule
  \endfirsthead
    \toprule
    #2 \\ \midrule
  \endhead
    \midrule
  \endfoot
    \bottomrule
  \endlastfoot
}
{\end{longtable}}

% --- Entorno TablaDatosFlexible (estilo antiguo, recuperado de config.tex) ---
%   Soportado para tablas migradas del trabajo previo. Usa longtable + booktabs
%   para consistencia con el estilo del rework (centrado, reglas profesionales,
%   cabecera repetible en cada página). Columnas L{} y C{} ya definidas arriba.
\newenvironment{TablaDatosFlexible}[2]{%
  \centering
  \setlength{\tabcolsep}{5pt}%
  \renewcommand{\arraystretch}{1.3}%
  \setlength{\LTleft}{0pt}%
  \setlength{\LTright}{0pt}%
  \begin{longtable}{#1}%
    \toprule
    #2 \\ \midrule
  \endfirsthead
    \toprule
    #2 \\ \midrule
  \endhead
    \midrule
  \endfoot
    \bottomrule
  \endlastfoot
}{%
  \end{longtable}%
  \par
}

\newcommand{\TablaFilaFlexible}[1]{#1 \\ \midrule}

% --- Caja reutilizable para resaltar resultados/interpretaciones ---
\newtcolorbox{CajaDestacada}[1][]{%
  colback=tablafila, colframe=azul!40, boxrule=0.6pt, arc=2pt,
  left=6pt, right=6pt, top=4pt, bottom=4pt, #1%
}

